\section{Related work}\label{sec:related}

\paragraph{Byzantine-robust aggregation.} From the Byzantine generals problem~\citep{lamport1982} to robust distributed learning, the dominant response to adversarial input is to \emph{remove} it. Krum~\citep{blanchard2017}, coordinate-wise median and trimmed mean~\citep{yin2018}, and Bulyan-style composition~\citep{guerraoui2018} are examples. These rules assume an honest majority and break at $f=1/2$. The original ``Liars Are Information'' study~\citep{das2026aip} observed that this ceiling belongs to \emph{discarding}, not to the problem. It proposed Adversarial Informativeness Pooling (\aip), which inverts channels whose coherence exceeds a calibrated honest ceiling. The same study then defeated \aip\ with an adversary that coordinates just below the ceiling, and proved that no coherence-measurable rule can separate such an adversary from honest disagreement. A derivative release~\citep{nabidnur2026bucket} added clean-task calibration, clone caps and causal online windows, with mixed results. \race\ keeps the premise (a liar can be read backwards) and replaces the coherence gate with channel estimation.

\paragraph{Learning from unreliable annotators.} Dawid and Skene~\citep{dawid1979} introduced per-annotator confusion matrices fitted by EM~\citep{dempster1977}. Later work added provably optimal spectral initialisation~\citep{zhang2016}, budget-optimal inference~\citep{karger2014}, joint learning of a classifier~\citep{raykar2010}, and unsupervised ensemble ranking~\citep{parisi2014}. Identifiability of such latent-class models holds only up to label permutation~\citep{kruskal1977,allman2009}. Crowdsourcing assumes a better-than-chance \emph{majority}, which is exactly what fails at $f\ge1/2$. Our contribution is to replace that assumption with a \emph{receiver} anchor (\cref{prop:switching}) and to handle replicated channels, which crowdsourcing rarely faces. The anchor resembles semi-supervised crowdsourcing, which seeds EM with gold questions or a trusted annotator, but it needs no gold: each receiver trusts only itself, every honest agent runs its own fit, and no party has to be trusted by the others. The ``surprisingly popular'' algorithm~\citep{prelec2017} also extracts truth from a minority. It needs a second elicitation (predicted votes) that an adversary could forge, whereas \race\ uses only the answers themselves.

\paragraph{Correlated LLM errors.} LLMs make correlated mistakes~\citep{kim2025correlated}. Judge panels are worth far fewer independent votes than their size suggests~\citep{kohli2026,verga2024juries,zheng2023judge}. Open-ended outputs are homogeneous across models~\citep{jiang2025hivemind}. Classical jury theorems degrade under correlated votes~\citep{ladha1992}. These findings motivate our error-conditioned clone detection (\cref{prop:clones}). We treat correlation as a quantity to estimate from errors, not from raw agreement.

\paragraph{Multi-agent LLM systems under attack.} Sampling-and-voting~\citep{wang2023sc} and multi-agent debate~\citep{du2024debate} improve reasoning when every agent is honest. Adversarial debaters can pull a group toward wrong answers~\citep{amayuelas2024}. Deceptive behaviour can persist in LLMs~\citep{hubinger2024sleeper,park2024deception}, and agents can be subverted by indirect prompt injection~\citep{greshake2023}. Our live experiment (\cref{sec:live}) measures the debate side directly. Exposing small honest models to a panel that is half liars barely changes their individual accuracy. On MMLU it removes information that receiver-local aggregation needs; on BoolQ, where the saboteurs carry almost no signal, it helps.
